\documentclass{article}

\usepackage{amsmath,amssymb}
\usepackage{xurl}
\usepackage[hidelinks]{hyperref}
\setlength{\emergencystretch}{2em}

% Shared commands for notes in docs/paper-gaps/.

\DeclareUnicodeCharacter{2080}{\ifmmode{}_{0}\else\textsubscript{0}\fi}
\DeclareUnicodeCharacter{2081}{\ifmmode{}_{1}\else\textsubscript{1}\fi}
\DeclareUnicodeCharacter{2082}{\ifmmode{}_{2}\else\textsubscript{2}\fi}
\DeclareUnicodeCharacter{2099}{\ifmmode{}_{n}\else\textsubscript{n}\fi}

\newcommand{\Lean}{\textsc{Lean}}
\ExplSyntaxOn
\NewDocumentCommand{\leanid}{m}{
  \ifmmode
    \textsf{\detokenize{#1}}
  \else
    \group_begin:
    \urlstyle{sf}
    \tl_set:Nn \l_tmpa_tl {#1}
    \tl_replace_all:Nnn \l_tmpa_tl {\_} {_}
    \exp_args:NV \path \l_tmpa_tl
    \group_end:
  \fi
}
\ExplSyntaxOff
\providecommand{\tr}{\operatorname{tr}}
\newcommand{\tnleanIssueBase}{https://github.com/LionSR/TNLean/issues/}
\newcommand{\tnleanPrBase}{https://github.com/LionSR/TNLean/pull/}
\newcommand{\tnleanDocBase}{https://sirui-lu.com/TNLean/docs/}
\newcommand{\ghissue}[1]{\href{\tnleanIssueBase#1}{GitHub issue \##1}}
\newcommand{\ghpr}[1]{\href{\tnleanPrBase#1}{PR \##1}}
\newcommand{\doclink}[1]{\href{\tnleanDocBase#1.html}{\path{#1}}}

% Dirac notation and the identity operator, as in blueprint/src/macros/common.tex.
% \providecommand keeps note-local definitions authoritative.
\usepackage{dsfont}
\providecommand{\ket}[1]{|#1\rangle}
\providecommand{\bra}[1]{\langle #1|}
\providecommand{\braket}[2]{\langle #1 | #2 \rangle}
\providecommand{\ketbra}[2]{|#1\rangle\!\langle #2|}
\providecommand{\Id}{\mathds{1}}
\providecommand{\one}{\mathds{1}}
\providecommand{\C}{\mathbb{C}}
\providecommand{\MN}[1]{M_{#1}(\C)}
\providecommand{\MD}{M_D(\C)}

% External referencing for paper sources.  The build helper writes sanitized
% auxiliary files under build/paper-aux/ and excludes duplicated source labels.
\usepackage{xr}

% Import a generated paper auxiliary file with a caller-supplied label prefix.
% The candidate paths support builds from docs/paper-gaps/, the repository root,
% and build/paper-aux/ itself.
\newcommand{\importpaperaux}[2]{%
  \IfFileExists{../../build/paper-aux/#2.aux}{%
    \externaldocument[#1]{../../build/paper-aux/#2}%
  }{%
    \IfFileExists{build/paper-aux/#2.aux}{%
      \externaldocument[#1]{build/paper-aux/#2}%
    }{%
      \IfFileExists{#2.aux}{%
        \externaldocument[#1]{#2}%
      }{}%
    }%
  }%
}

% General external reference macros
\newcommand{\paperref}[2]{\ref{#1-#2}}
\newcommand{\papereqref}[2]{(\ref{#1-#2})}

% CPSV16 source (1606.00608 / MPDO-22-12-17-2.tex) external document and shortcuts
\importpaperaux{cpsv16-}{1606.00608/MPDO-22-12-17-2-tnlean}
\newcommand{\cpsvref}[1]{\paperref{cpsv16}{#1}}
\newcommand{\cpsveqref}[1]{\papereqref{cpsv16}{#1}}

% Verdict marker: \gapnote{<kind>}{<status>}. Kind is the policy
% classification (clarification, local-correction, scope-restriction,
% unfaithful, false-source, open-gap); status is open, wip (elimination
% underway), resolved, or historical. Machine-read by the site index;
% prints a verdict line.
\newcommand{\gapnote}[2]{\par\noindent\textbf{Verdict:} #1 (#2).\par}


\title{Wolf Proposition~6.2 --- Trivial Jordan Blocks for Peripheral Eigenvalues}
\author{}
\date{\today}

\begin{document}
\maketitle

\section{The assertion}

Wolf, \emph{Quantum Channels \& Operations}, Proposition~6.2:
``Let $T:\mathcal M_d(\mathbb C)\to\mathcal M_d(\mathbb C)$ be a
trace-preserving (or unital) positive linear map.  If $\lambda$ is an
eigenvalue of $T$ with $|\lambda|=1$, then its geometric multiplicity
equals its algebraic multiplicity, i.e., all Jordan blocks for $\lambda$
are one-dimensional.''

The proof observes that the iterates $T^n$ are uniformly bounded in
operator norm (positivity and trace preservation keep every forward orbit
bounded) and that powers of a nontrivial Jordan block grow without bound
when the eigenvalue has modulus~$1$, giving a contradiction.

\section{The formalization}

The bounded-orbit lemma is already formalized:
\leanid{IsPositiveMap.hasBoundedOrbits\_of\_tracePreserving} in
\doclink{TNLean/Channel/FixedPoint/MeanErgodicProjection} proves
that every forward orbit $\{T^n X\}_{n\in\mathbb N}$ is bounded in norm
for a positive trace-preserving $T$.

The binomial-expansion ingredient has not been formalized.  A rank-$2$
generalized eigenvector $(T-\lambda)X=Y\neq0$ with $(T-\lambda)Y=0$
would give, by \leanid{Commute.add\_pow} applied to the decomposition
$T=\lambda I+(T-\lambda I)$,
\[
    T^n X = \lambda^n X + n\lambda^{n-1}Y,
\]
whose norm grows linearly with $n$, contradicting bounded orbits.

The missing work consists in (i) packaging \leanid{Commute.add\_pow}
for the semiring of linear endomorphisms to obtain this identity,
(ii) extracting a norm bound $\|T^n X\|\ge n\|Y\|-\|X\|$ from it, and
(iii) applying the bounded-orbit lemma to reach a contradiction.
The full statement for higher-rank generalized eigenvectors
($(T-\lambda)^k X=0$ with $k>1$) would additionally require the
binomial term $\binom{n}{k-1}$ growing like $n^{k-1}$.

\section{Verdict}

This is an \emph{open gap}.  The bounded-orbit lemma supplies the analytic
hypothesis; the algebraic step (binomial expansion for the power of a linear
endomorphism with nilpotent perturbation) is mathematically clear but not
yet formalized.  The rank-$2$ case suffices for the peripheral-spectrum
applications in Theorem~6.6 of the source.\footnote{Tracking: \ghissue{3984}.}

\bibliographystyle{alpha}
\bibliography{references}

\end{document}
